\documentclass[10pt,letterpaper]{article}
\usepackage[utf8]{inputenc}
\usepackage[intlimits]{amsmath}
\usepackage{amsfonts}
\usepackage{amssymb}
\usepackage{ragged2e}
\usepackage[letterpaper, margin=1in]{geometry}
\usepackage{graphicx}
\usepackage{cancel}
\usepackage{mathtools}
\usepackage{tabularx}
\usepackage{arydshln}
\usepackage{tensor}
\usepackage{array}
\usepackage{xcolor}
\usepackage[boxed]{algorithm}
\usepackage[noend]{algpseudocode}
\usepackage{listings}
\usepackage{textcomp}
% \usepackage[pdf,tmpdir,singlefile]{graphviz}
\usepackage{mathrsfs}
\usepackage{bbm}
\usepackage{tikz}
\usepackage{tikz-cd}
\usepackage{enumitem}
\usepackage{arydshln}
\usepackage{relsize}
\usepackage{multicol}
\usepackage{scalerel}

\usetikzlibrary{bayesnet}
\setlist{noitemsep}

%%%%%%%%%%%%%%%%%%%%%%%%%%%%%
% Formatting commands
%%%%%%%%%%%%%%%%%%%%%%%%%%%%%
\newcommand{\n}{\hfill\break}
\newcommand{\up}{\vspace{-\baselineskip}}
\newcommand{\hangblock}[2]{\par\noindent\settowidth{\hangindent}{\textbf{#1: }}\textbf{#1: }\!\!\!#2}
\newcommand{\lemma}[1]{\hangblock{Lemma}{#1}}
\newcommand{\defn}[1]{\hangblock{Defn}{#1}}
\newcommand{\thm}[1]{\hangblock{Thm}{#1}}
\newcommand{\cor}[1]{\hangblock{Cor}{#1}}
\newcommand{\prop}[1]{\hangblock{Prop}{#1}}
\newcommand{\ex}[1]{\hangblock{Ex}{#1}}
\newcommand{\exer}[1]{\hangblock{Exer}{#1}}
\newcommand{\fact}[1]{\hangblock{Fact}{#1}}
\newcommand{\remark}[1]{\hangblock{Remark}{#1}}
\newcommand{\proven}{\;$\square$\n}
\newcommand{\problem}[1]{\par\noindent{#1}\n}
\newcommand{\problempart}[2]{\par\noindent\indent{}\settowidth{\hangindent}{\textbf{(#1)} \indent{}}\textbf{(#1) }\!\!\!#2\n}
\newcommand{\ptxt}[1]{\textrm{\textnormal{#1}}}
\newcommand{\inlineeq}[1]{\centerline{$\displaystyle #1$}}
\newcommand{\pageline}{\up\par\noindent\rule{\textwidth}{0.1pt}}

%%%%%%%%%%%%%%%%%%%%%%%%%%%%%
% Math commands
%%%%%%%%%%%%%%%%%%%%%%%%%%%%%
% Set Theory
\newcommand{\card}[1]{\left|#1\right|}
\newcommand{\set}[1]{\left\{#1\right\}}
\newcommand{\setmid}{\;\middle|\;}
\newcommand{\ps}[1]{\mathcal{P}\left(#1\right)}
\newcommand{\pfinite}[1]{\mathcal{P}^{\ptxt{finite}}\left(#1\right)}
\newcommand{\naturals}{\mathbb{N}}
\newcommand{\N}{\naturals}
\newcommand{\integers}{\mathbb{Z}}
\newcommand{\Z}{\integers}
\newcommand{\rationals}{\mathbb{Q}}
\newcommand{\Q}{\rationals}
\newcommand{\reals}{\mathbb{R}}
\newcommand{\R}{\reals}
\newcommand{\complex}{\mathbb{C}}
\newcommand{\C}{\complex}
\newcommand{\halfPlane}{\mathbb{H}}
\let\H\relax
\newcommand{\H}{\halfPlane}
\newcommand{\comp}{^{\complement}}
\DeclareMathOperator{\Hom}{Hom}
\newcommand{\Ind}{\mathbbm{1}}
\newcommand{\cut}{\setminus}
\DeclareMathOperator{\elem}{elem}

% Graph Theory
\let\deg\relax
\DeclareMathOperator{\deg}{deg}
\newcommand{\degp}{\ptxt{deg}^{+}}
\newcommand{\degn}{\ptxt{deg}^{-}}
\newcommand{\precdot}{\mathrel{\ooalign{$\prec$\cr\hidewidth\hbox{$\cdot\mkern0.5mu$}\cr}}}
\newcommand{\succdot}{\mathrel{\ooalign{$\cdot\mkern0.5mu$\cr\hidewidth\hbox{$\succ$}\cr\phantom{$\succ$}}}}
\DeclareMathOperator{\cl}{cl}
\DeclareMathOperator{\affdim}{affdim}

% Probability
\newcommand{\parSymbol}{\P}
\newcommand{\Prob}{\mathbb{P}}
\renewcommand{\P}{\Prob}
\newcommand{\Avg}{\mathbb{E}}
\newcommand{\E}{\Avg}
\DeclareMathOperator{\Var}{Var}
\DeclareMathOperator{\cov}{cov}
\DeclareMathOperator{\Unif}{Unif}
\DeclareMathOperator{\Binom}{Binom}
\newcommand{\CI}{\mathrel{\text{\scalebox{1.07}{$\perp\mkern-10mu\perp$}}}}

% Standard Math
\newcommand{\inv}{^{-1}}
\newcommand{\abs}[1]{\left|#1\right|}
\newcommand{\ceil}[1]{\left\lceil{}#1\right\rceil{}}
\newcommand{\floor}[1]{\left\lfloor{}#1\right\rfloor{}}
\newcommand{\conj}[1]{\overline{#1}}
\newcommand{\of}{\circ}
\newcommand{\tri}{\triangle}
\newcommand{\inj}{\hookrightarrow}
\newcommand{\surj}{\twoheadrightarrow}
\newcommand{\ndiv}{\nmid}
\renewcommand{\epsilon}{\varepsilon}
\newcommand{\divides}{\mid}
\newcommand{\ndivides}{\nmid}
\DeclareMathOperator{\lcm}{lcm}
\DeclareMathOperator{\sgn}{sgn}
\newcommand{\map}[4]{\!\!\!\begin{array}[t]{rcl}#1 & \!\!\!\!\to & \!\!\!\!#2\\ #3 & \!\!\!\!\mapsto & \!\!\!\!#4\end{array}}
\newcommand{\bigsum}[2]{\smashoperator[lr]{\sum_{\scalebox{#1}{$#2$}}}}

% Linear Algebra
\newcommand{\Id}{\textrm{\textnormal{Id}}}
\newcommand{\im}{\textrm{\textnormal{im}}}
\newcommand{\norm}[1]{\abs{\abs{#1}}}
\newcommand{\tpose}{^{T}\!}
\newcommand{\iprod}[1]{\left<#1\right>}
\newcommand{\giprod}{\iprod{\;\,,\;}}
\DeclareMathOperator{\tr}{tr}
\DeclareMathOperator{\trace}{tr}
\newcommand{\chgBasMat}[3]{\!\!\tensor*[_{#1}]{\left[#2\right]}{_{#3}}}
\newcommand{\vecBas}[2]{\tensor*[]{\left[#1\right]}{_{#2}}}
\DeclareMathOperator{\GL}{GL}
\DeclareMathOperator{\Mat}{Mat}
\DeclareMathOperator{\vspan}{span}
\DeclareMathOperator{\rank}{rank}
\newcommand{\V}[1]{\vec{#1}}
\DeclareMathOperator{\proj}{proj}
\DeclareMathOperator{\compProj}{comp}
\DeclareMathOperator{\row}{row}
\newcommand{\smallPMatrix}[1]{\paren{\begin{smallmatrix}#1\end{smallmatrix}}}
\newcommand{\smallBMatrix}[1]{\brack{\begin{smallmatrix}#1\end{smallmatrix}}}
\newcommand{\pmat}[1]{\begin{pmatrix}#1\end{pmatrix}}
\newcommand{\bmat}[1]{\begin{bmatrix}#1\end{bmatrix}}

% Multilinear Algebra
\newcommand{\Lsym}{\L}
\let\L\relax
\DeclareMathOperator{\L}{\mathscr{L}}
\DeclareMathOperator{\A}{\mathcal{A}}
\DeclareMathOperator{\Alt}{Alt}
\DeclareMathOperator{\Sym}{Sym}
\newcommand{\ot}{\otimes}
\newcommand{\ox}{\otimes}
\DeclareMathOperator{\asc}{asc}
\DeclareMathOperator{\asSet}{set}
\DeclareMathOperator{\sort}{sort}
\DeclareMathOperator{\ringA}{\mathring{A}}

% Topology
\newcommand{\closure}[1]{\overline{#1}}
\newcommand{\uball}{\mathcal{U}}
\DeclareMathOperator{\Int}{Int}
\DeclareMathOperator{\Ext}{Ext}
\DeclareMathOperator{\Bd}{Bd}
\DeclareMathOperator{\rInt}{rInt}
\DeclareMathOperator{\ch}{ch}
\DeclareMathOperator{\ah}{ah}
\newcommand{\LargerTau}{\mathlarger{\mathlarger{\mathlarger{\mathlarger{\tau}}}}}
\newcommand{\Tau}{\mathcal{T}}

% Analysis
\DeclareMathOperator{\Graph}{Graph}
\DeclareMathOperator{\epi}{epi}
\DeclareMathOperator{\hypo}{hypo}
\DeclareMathOperator{\supp}{supp}
\newcommand{\lint}[2]{\underset{#1}{\overset{#2}{{\color{black}\underline{{\color{white}\overline{{\color{black}\int}}\color{black}}}}}}}
\newcommand{\uint}[2]{\underset{#1}{\overset{#2}{{\color{white}\underline{{\color{black}\overline{{\color{black}\int}}\color{black}}}}}}}
\newcommand{\alignint}[2]{\underset{#1}{\overset{#2}{{\color{white}\underline{{\color{white}\overline{{\color{black}\int}}\color{black}}}}}}}
\newcommand{\extint}{\ptxt{ext}\int}
\newcommand{\extalignint}[2]{\ptxt{ext}\alignint{#1}{#2}}
\newcommand{\conv}{\ast}
\newcommand{\pd}[2]{\frac{\partial{}#1}{\partial{}#2}}
\newcommand{\del}{\nabla}
\DeclareMathOperator{\grad}{grad}
\DeclareMathOperator{\curl}{curl}
\let\div\relax
\DeclareMathOperator{\div}{div}
\DeclareMathOperator{\vol}{vol}

% Complex Analysis
\let\Re\relax
\DeclareMathOperator{\Re}{Re}
\let\Im\relax
\DeclareMathOperator{\Im}{Im}
\DeclareMathOperator{\Res}{Res}

% Abstract Algebra
\DeclareMathOperator{\ord}{ord}
\newcommand{\generated}[1]{\left<#1\right>}
\newcommand{\cycle}[1]{\smallPMatrix{#1}}
\newcommand{\id}{\ptxt{id}}
\newcommand{\iso}{\cong}
\DeclareMathOperator{\Aut}{Aut}
\DeclareMathOperator{\SL}{SL}
\DeclareMathOperator{\op}{op}
\newcommand{\isom}[4]{\!\!\!\begin{array}[t]{rcl}#1 & \!\!\!\!\overset{\sim}{\to} & \!\!\!\!#2\\ #3 & \!\!\!\!\mapsto & \!\!\!\!#4\end{array}}
\newcommand{\F}{\mathbb{F}}
\newcommand{\acts}{\;\reflectbox{\rotatebox[origin=c]{-90}{$\circlearrowright$}}\;}
\newcommand{\disjunion}{\mathrel{\text{\scalebox{1.07}{$\perp\mkern-10mu\perp$}}}}
\DeclareMathOperator{\SO}{SO}
\DeclareMathOperator{\stab}{stab}
\DeclareMathOperator{\nullity}{nullity}
\DeclareMathOperator{\Perm}{Perm}
\DeclareMathOperator{\nsubgp}{\vartriangleleft}
\DeclareMathOperator{\notnsubgp}{\ntriangleleft}
\newcommand{\presentation}[2]{\left<#1\;\middle|\;#2\right>}

% Convex Optimization
\newcommand{\sectionSymbol}{\S}
\let\S\relax
\newcommand{\S}{\mathbb{S}}
\DeclareMathOperator{\dist}{dist}
\DeclareMathOperator{\dom}{dom}
\DeclareMathOperator{\diag}{diag}
\DeclareMathOperator{\ones}{\mathbbm{1}}
\newcommand{\minimizeOver}[3]{\begin{array}{rl}\underset{#1}{\ptxt{minimize}} & #2\\ \ptxt{subject to} & #3\end{array}}
\newcommand{\maximizeOver}[3]{\begin{array}{rl}\underset{#1}{\ptxt{maximize}} & #2\\ \ptxt{subject to} & #3\end{array}}
\newcommand{\minimizationProblem}[2]{\minimizeOver{}{#1}{#2}}
\newcommand{\maximizationProblem}[2]{\maximizeOver{}{#1}{#2}}
\newcommand{\minimizeOverUnconstrained}[2]{\begin{array}{rl}\underset{#1}{\ptxt{minimize}} & #2\end{array}}
\newcommand{\maximizeOverUnconstrained}[2]{\begin{array}{rl}\underset{#1}{\ptxt{maximize}} & #2\end{array}}
\newcommand{\minimizationUnconstrained}[1]{\minimizeOverUnconstrained{}{#1}}
\newcommand{\maximizationUnconstrained}[1]{\maximizeOverUnconstrained{}{#1}}

% Proofs
\newcommand{\st}{s.t.}
\newcommand{\unique}{!}
\newcommand{\iffdef}{\overset{\ptxt{def}}{\Leftrightarrow}}
\newcommand{\eqdef}{\overset{\ptxt{def}}{=}}
\newcommand{\eqVertical}{\rotatebox[origin=c]{90}{=}}
\newcommand{\mapsfrom}{\mathrel{\reflectbox{\ensuremath{\mapsto}}}}
\newcommand{\mapsdown}{\rotatebox[origin=c]{-90}{$\mapsto$}\mkern2mu}
\newcommand{\mapsup}{\rotatebox[origin=c]{90}{$\mapsto$}\mkern2mu}
\newcommand{\from}{\!\mathrel{\reflectbox{\ensuremath{\to}}}}

% Brackets
\newcommand{\paren}[1]{\left(#1\right)}
\renewcommand{\brack}[1]{\left[#1\right]}
\renewcommand{\brace}[1]{\left\{#1\right\}}
\newcommand{\ang}[1]{\left<#1\right>}

% Algorithms
\algrenewcommand{\algorithmiccomment}[1]{\hskip 1em \texttt{// #1}}
\algrenewcommand\algorithmicrequire{\textbf{Input:}}
\algrenewcommand\algorithmicensure{\textbf{Output:}}
\newcommand{\algP}{\ptxt{\textbf{P}}}
\newcommand{\algNP}{\ptxt{\textbf{NP}}}
\newcommand{\algNPC}{\ptxt{\textbf{NP-Complete}}}
\newcommand{\algNPH}{\ptxt{\textbf{NP-Hard}}}
\newcommand{\algEXP}{\ptxt{\textbf{EXP}}}

%%%%%%%%%%%%%%%%%%%%%%%%%%%%%
% Other commands
%%%%%%%%%%%%%%%%%%%%%%%%%%%%%
\newcommand{\flag}[1]{\textbf{\textcolor{red}{#1}}}
\newcommand{\uSym}{\u}
\let\u\relax
\newcommand{\u}[1]{\underline{#1}}
\newcommand{\bSym}{\b}
\let\b\relax
\newcommand{\b}[1]{\textbf{#1}}
\newcommand{\iSym}{\i}
\let\i\relax
\newcommand{\i}[1]{\textit{#1}}

%%%%%%%%%%%%%%%%%%%%%%%%%%%%%%%%%%%%%%%
% Make l's curvy in math environments %
%%%%%%%%%%%%%%%%%%%%%%%%%%%%%%%%%%%%%%%
\mathcode`l="8000
\begingroup
\makeatletter
\lccode`\~=`\l
\DeclareMathSymbol{\lsb@l}{\mathalpha}{letters}{`l}
\lowercase{\gdef~{\ifnum\the\mathgroup=\m@ne \ell \else \lsb@l \fi}}%
\endgroup

%%%%%%%%%%%%%%%%%%%%%%%%%
% Fix \vdots and \ddots %
%%%%%%%%%%%%%%%%%%%%%%%%%
\usepackage{letltxmacro}
\LetLtxMacro\orgvdots\vdots
\LetLtxMacro\orgddots\ddots

\makeatletter
\DeclareRobustCommand\vdots{%
	\mathpalette\@vdots{}%
}
\newcommand*{\@vdots}[2]{%
	% #1: math style
	% #2: unused
	\sbox0{$#1\cdotp\cdotp\cdotp\m@th$}%
	\sbox2{$#1.\m@th$}%
	\vbox{%
		\dimen@=\wd0 %
		\advance\dimen@ -3\ht2 %
		\kern.5\dimen@
		% remove side bearings
		\dimen@=\wd2 %
		\advance\dimen@ -\ht2 %
		\dimen2=\wd0 %
		\advance\dimen2 -\dimen@
		\vbox to \dimen2{%
			\offinterlineskip
			\copy2 \vfill\copy2 \vfill\copy2 %
		}%
	}%
}
\DeclareRobustCommand\ddots{%
	\mathinner{%
		\mathpalette\@ddots{}%
		\mkern\thinmuskip
	}%
}
\newcommand*{\@ddots}[2]{%
	% #1: math style
	% #2: unused
	\sbox0{$#1\cdotp\cdotp\cdotp\m@th$}%
	\sbox2{$#1.\m@th$}%
	\vbox{%
		\dimen@=\wd0 %
		\advance\dimen@ -3\ht2 %
		\kern.5\dimen@
		% remove side bearings
		\dimen@=\wd2 %
		\advance\dimen@ -\ht2 %
		\dimen2=\wd0 %
		\advance\dimen2 -\dimen@
		\vbox to \dimen2{%
			\offinterlineskip
			\hbox{$#1\mathpunct{.}\m@th$}%
			\vfill
			\hbox{$#1\mathpunct{\kern\wd2}\mathpunct{.}\m@th$}%
			\vfill
			\hbox{$#1\mathpunct{\kern\wd2}\mathpunct{\kern\wd2}\mathpunct{.}\m@th$}%
		}%
	}%
}
\makeatother

\newcommand{\B}{
	\begin{tikzpicture}
	\filldraw [fill=red, draw=black] (0, 0) rectangle (0.37, 0.45);
	\draw [line width=0.5mm, white ] (0.1,0.08) -- (0.1,0.38);
	\draw[line width=0.5mm, white ] (0.1, 0.35) .. controls (0.2, 0.35) and (0.4, 0.2625) .. (0.1, 0.225);
	\draw[line width=0.5mm, white ] (0.1, 0.225) .. controls (0.2, 0.225) and (0.4, 0.1625) .. (0.1, 0.1);
	\end{tikzpicture}
}

\author{Professor Andrew Snowden\\ \small\i{Transcribed by Thomas Cohn}}
\title{Math 493 Lecture 20}
\date{11/13/19} % Can also use \today

\begin{document}
\maketitle
\setlength\RaggedRightParindent{\parindent}
\RaggedRight

\par\noindent
Recall from last time: We have $F$, a field not of characteristic $2$, $V$, a finite dimensional $F$-vector space, and $\giprod$, a symmetric bilinear form on $V$ ($\giprod:V\times{}V\to{}F$).\n

\prop{
	(From last time) there is an orgthogonal basis for $V$, i.e., a basis $(e_{1},\ldots,e_{n})$ \st{} $\iprod{e_{i},e_{j}}=0$ for $i\ne{}j$, if and only if $V\cong[a_{1},\ldots,a_{n}]$ for some choice of $a_{1},\ldots,a_{n}\in{}F$, for which $a_{i}=\iprod{e_{i},e_{i}}$.\n
}

\par\noindent
Notation: given $a_{1},\ldots,a_{n}\in{}F$, $[a_{1},\ldots,a_{n}]$ is the quadratic space with vector space $F^{n}$ and the form is
\[
	\iprod{e_{i},e_{j}}=\left\{\begin{array}{ll}0 & i\ne{}j\\ a_{i} & i=j\end{array}\right.
\]

\par\noindent
For $[a_{1},\ldots,a_{n}]$ in the standard basis, the matrix of the form is the diagonal matrix
\[
	\begin{bmatrix}
		a_{1} & & 0\\
		& \ddots &\\
		0 & & a_{n}
	\end{bmatrix}
\]

\par\noindent
Recall: The kernel of a quadratic space $V$ is $\set{v\in{}V\mid\iprod{v,w}=0,\forall{}w\in{}V}$.\n

\defn{
	We say $V$ is \textbf{non-degenerate} if $\ker{}V=\set{0}$.\n
}

\prop{
	Say $V=[a_{1},\ldots,a_{n}]$. The following are equivalent:
	\begin{enumerate}[leftmargin=4\parindent, label=(\arabic*)]
		\item $V$ is non-degenerate
		\item $a_{i}\ne{}0$, $\forall{}i$
		\item The discriminant is nonzero
	\end{enumerate}\up\n
	Proof: (2) $\Leftrightarrow$ (3) is clear because the discriminant is $a_{1}\cdots{}a_{n}$.\n
	If some $a_{i}=0$, then $\iprod{e_{i},e_{j}}=0$, $\forall{}j$. So $e_{i}\in\ker{}V$, as $\iprod{e_{i},\sum_{j=1}^{n}\alpha_{j}e_{j}}=\sum_{j=1}^{n}\alpha_{j}\iprod{e_{i},e_{j}}=0$. So $V$ is degenerate.\n
	Say all $a_{i}$ are nonzero. Then let $v=\alpha_{1}e_{1}+\cdots+\alpha_{n}e_{n}$ be a nonzero element of $V$. So $\exists\alpha_{i}\ne{}0$. So $\iprod{v,e_{i}}=\alpha_{i}\iprod{e_{i},e_{i}}=\alpha_{i}a_{i}\ne{}0$. So $v\not\in\ker{}V$, so $\ker{}V=\set{0}$.\proven
}

\cor{
	Every $V$ is isomorphic to $U\perp{}W$, where the form is identically $0$ on $U$ and $W$ is non-degenerate.\n
	Proof: Write $V=[a_{1},\ldots,a_{n}]$. Say $a_{1},\ldots,a_{m}=0,a_{m+1},\ldots,a_{n}\ne{}0$. Then let $U=[a_{1},\ldots,a_{m}]$ and $W=[a_{m+1},\ldots,a_{n}]$.\proven
}

\par\noindent
This is great! It allows us to basically always work with non-degenerate forms.\n

\par\noindent
Observe that for any field $F$, we have $[a]\cong[ab^{2}]$, for any $b\in\F^{\times}$.\n
Proof: Let $V=[a]$ for basis $e$ with $\iprod{e,e}=a$. Let $W=[ab^{2}]$ for basis $f$ with $\iprod{f,f}=ab^{2}$. Then define a linear map $T:V\to{}W$ where $e\mapsto\frac{1}{b}f$. We have
\[
	\iprod{T(e),T(e)}=\iprod{\frac{1}{b}f,\frac{1}{b}f}=\paren{\frac{1}{b}}^{2}\iprod{f,f}=\paren{\frac{1}{b}}^{2}ab^{2}=a=\iprod{e,e}
\]
so $T$ is an isometry. Thus, $V\cong{}W$. In fact, we have $[a_{1},\ldots,a_{n}]\cong[a_{1}b_{1}^{2},\ldots,a_{n}b_{n}^{2}]$ for any $b_{1},\ldots,b_{n}\in{}F^{\times}$.\proven

\prop{
	Two non-degenerate quadratic spaces over $\C$ are isometric iff they have the same dimension.\n
	Proof: The same dimension requirement is obviously necessary.\n
	By our previous observation, $[a_{1},\ldots,a_{n}]\cong[1,\ldots,1]$ if $a_{1},\ldots,a_{n}$ are all nonzero (we can select $b_{i}=\frac{1}{\sqrt{a_{i}}}$). In particular, $[a_{1},\ldots,a_{n}]\cong[a'_{1},\ldots,a'_{n}]$ if both are non-degenerate.\proven
}

\par\noindent
Over $\R$, if $a>0$, then $[a]\cong[1]$, and if $a<0$, then $[a]\cong[-1]$. (In both cases, take $b=1/\sqrt{\abs{a}}$). In general,
\[
	[a_{1},\ldots,a_{n}]\cong[\underbrace{1,\ldots,1}_{r},\underbrace{-1,\ldots,-1}_{n-r}]
\]

\thm{
	(Sylvester's Law of Inertia) The number $r$ is well-defined, i.e.,
	\[
		[\underbrace{1,\ldots,1}_{p},-1,\ldots,-1]\cong[\underbrace{1,\ldots,1}_{q},-1,\ldots,-1]\quad\Rightarrow\quad{}p=q
	\]
	Proof: Let $V$ be a non-degenerate quadratic space. Say $e_{1},\ldots,e_{n}$ and $f_{1},\ldots,f_{n}$ are orgthogonal bases \st{}
	\[
		\iprod{e_{i},e_{i}}=\left\{\begin{array}{ll}
			1 & 1\le{}i\le{}p\\
			-1 & p<i\le{}n
		\end{array}\right.
		\quad\ptxt{and}\quad
		\iprod{f_{i},f_{i}}=\left\{\begin{array}{ll}
			1 & 1\le{}i\le{}q\\
			-1 & q<i\le{}n
		\end{array}\right.
	\]
	Let $U=\vspan(e_{1},\ldots,e_{p})$ and $W=\vspan(f_{q+1},\ldots,f_{n})$. We claim $U\cap{}W=\set{0}$.\n
	Let $v\in{}U\cap{}W$. Then
	\begin{align*}
		& v\in{}V{}\quad\Rightarrow\quad{}v=\alpha_{1}e_{1}+\cdots+\alpha_{p}e_{p},\alpha_{i}\in\R\qquad\Rightarrow\quad\iprod{v,v}=\sum_{i=1}^{p}\alpha_{i}^{2}\ge{}0\\
		& v\in{}\mathrlap{W}\phantom{V}\quad\Rightarrow\quad{}\mathrlap{v=\beta_{q+1}f_{1}+\cdots+\beta_{n}f_{n},\beta_{i}\in\R}\phantom{v=\alpha_{1}e_{1}+\cdots+\alpha_{p}e_{p},\alpha_{i}\in\R}\qquad\Rightarrow\quad\iprod{v,v}=\smashoperator{\sum_{i=q+1}^{n}}\beta_{i}^{2}\le{}0
	\end{align*}
	So $\iprod{v,v}=0$, so $v=0$.\n
	Thus, we have $U\cap{}W=\set{0}$, so $\dim{}U+\dim{}W\le{}\dim{}V$. Thus, $p+(n-q)\le{}n$, so $p\le{}q$. We can now repeat our argument in the opposite direction, to obtain $q\le{}p$, so we have $q=p$.\proven
}

\par\noindent
In summary, if $V$ is a non-degenerate quadratic space over $\R$ of dimension $n$, $\exists\unique{}r,s$ \st{} $r+s=n$ and $V\cong[\underbrace{1,\ldots,1}_{r},\underbrace{-1,\ldots,-1}_{s}]$.\n

\defn{
	$(r,s)$ is the \textbf{signature} of $V$.\n
}

\par\noindent
Let $F=\F_{p}$ ($p$ odd, i.e., $p\ne{}2$).\n

\prop{
	$F^{\times}/(F^{\times})^{2}\cong\Z/2\Z$.\n
	Proof: $\F^{\times}$ is cyclic of even order.\proven
	Proof 2: Let $f:F^{\times}\to{}F^{\times}$ take $x\mapsto{}x^{2}$.\n
	This is a group homomorphism -- $f(xy)=(xy)^{2}=x^{2}y^{2}=f(x)f(y)$.\n
	$\im(f)=(F^{\times})^{2}$. $\ker(f)=\set{x\in{}F^{\times}\mid{}x^{2}=1}$. Well,
	\begin{align*}
		x^{2}=1 & \Leftrightarrow{}x^{2}-1=0\\
		& \Leftrightarrow{}(x+1)(x-1)=0\\
		& \Leftrightarrow{}x+1=0\;\ptxt{or}\;{}x-1=0\\
		& \Leftrightarrow{}x=-1\;\ptxt{or}\;{}x=1
	\end{align*}
	So by the First Isomorphism Theorem, $\#(F^{\times})^{2}=\#\im(f)=\frac{\#{}F^{\times}}{\#\ker(f)}=\frac{\#F^{\times}}{2}$.\n
	Thus, $\#F^{\times}/(F^{\times})^{2}=2$.\proven
}

\par\noindent
In summary, $\exists{}a\in{}F^{\times}$ \st{} every element of $F^{\times}$ has the form $b^{2}$ or $ab^{2}$, for some $b\in{}F^{\times}$. This $a$ is a non-square, and is not unique. Sometimes, we can have $a=-1$; other times, we cannot. In fact, $-1$ is a square iff $p\equiv{}1\pmod{4}$.\n

\ex{
	Let $p=43$. Then $p\equiv{}3\pmod{4}$, so $-1$ is not a square, so $a=-1$ is allowed. So every element of $\F_{43}^{\times}$ has the form $\pm{}b^{2}$ for some $b\in\F_{43}^{\times}$.\n
}

\ex{
	Let $p=41$. Then $p\equiv{}1\pmod{4}$, so $-1$ is a square, so we can't use $a=-1$. $2$ is also bad, as\n
	$2=17^{2}\bmod{41}$.\n
}

\exer{
	Find some value for $a$ for $p=41$.\n
}

\par\noindent
Fix $\varepsilon\in{}F^{\times}$ not a square. Just as in the real case, for any $a_{1},\ldots,a_{n}\in{}F^{\times}$, we have an isomorphism $[a_{1},\ldots,a_{n}]\cong[\underbrace{1,\ldots,1}_{r},\underbrace{\varepsilon,\ldots,\varepsilon}_{n-r}]$. But $r$ is \textbf{not} well-defined -- only $r\pmod{2}$ is, because $[\varepsilon,\varepsilon]\cong[1,1]$.\n

\par\noindent
The key point of all of this is every element of $F$ is a sum of $2$ squares.

\end{document}